\documentclass[11pt]{article}

\usepackage{amsmath,amssymb,amsthm}
\usepackage{graphicx}
\usepackage{booktabs}
\usepackage{hyperref}
\usepackage{natbib}
\usepackage[margin=1.2in]{geometry}
\usepackage{microtype}
\usepackage{xcolor}
\usepackage{caption}
\usepackage{subcaption}

\hypersetup{colorlinks=true, linkcolor=blue!60!black, citecolor=blue!60!black, urlcolor=blue!60!black}

\newtheorem{proposition}{Proposition}
\newcommand{\E}{\mathbb{E}}
\newcommand{\R}{\mathbb{R}}
\newcommand{\Var}{\mathrm{Var}}
\newcommand{\N}{\mathcal{N}}

% ── Title block ──────────────────────────────────────────────────────────────
\title{%
  Probabilistic College Basketball Ratings:\\
  A Ridge Prior Is Worth Three Weeks of Data
}
\author{Aaron Danielson}
\date{\today}

% ─────────────────────────────────────────────────────────────────────────────
\begin{document}
\maketitle

\begin{abstract}
We develop a probabilistic framework for college basketball team ratings that
produces a full posterior distribution over team quality, rather than a point
estimate.  Two models are derived and compared: a multiplicative fixed-point
system modeled after KenPom (Model~1) and a ridge regression model with an
exact Gaussian posterior (Model~2).  Both are validated under a strict
\emph{one-step-ahead} protocol on the 2025--26 NCAA season: for each week, the
model is fit on all prior games and evaluated on that week's games exclusively.
Model~2 achieves a season-wide RMSE of~14.14 pts/100 possessions versus~15.10
for Model~1.  The gap is largest early in the season, where Model~1 is nearly
underdetermined: in week~1, Model~1 RMSE exceeds~22 while Model~2 opens at
14.6 --- empirically confirming that the ridge prior provides the equivalent of
approximately three weeks of additional data.  Model~2's posterior predictive
intervals are empirically well-calibrated: a 90\% interval contains roughly
90\% of actual outcomes across all coverage levels.  The small systematic
negative bias observed in both models is a direct consequence of within-season
team improvement, pointing naturally at a dynamic-coefficient extension.
\end{abstract}

% ─────────────────────────────────────────────────────────────────────────────
\section{Introduction}

College basketball rating systems --- KenPom \citep{pomeroy2024}, BartTorvik,
ESPN's BPI --- universally report point estimates of team quality.  A team is
assigned an adjusted offensive efficiency, a defensive efficiency, and a tempo,
and these numbers are used for ranking, seeding, and matchup analysis.

This approach has a structural limitation: point estimates give no indication of
how confident the model is.  A team that has played three games has the same
representation as a team that has played thirty.  Early in the season, when
teams have played five to ten games each and the schedule graph is sparse, the
estimates are largely determined by prior assumptions that are never made
explicit.

We make those assumptions explicit.  Specifically, we show that:

\begin{enumerate}
  \item A Gaussian prior over team effects (the implicit assumption of ridge
    regression) produces well-posed inference from the first game of the
    season, with quantifiable uncertainty.
  \item This prior is worth approximately three weeks of additional data in
    early-season forecast accuracy, relative to the fixed-point system that has
    no explicit regularization.
  \item The resulting posterior predictive intervals are empirically
    well-calibrated across all coverage levels, enabling honest probabilistic
    statements about matchup outcomes.
  \item The systematic underprediction bias observed in both models is
    interpretable as a within-season team improvement effect, and points
    directly at a state-space extension.
\end{enumerate}

The framework is built on a common data structure --- offensive efficiency
per 100 possessions, one observation per team per game --- and produces a
KenPom-compatible summary (AdjO, AdjD, AdjPace) as a \emph{derived} quantity
from the posterior, rather than a primitive estimate.

% ─────────────────────────────────────────────────────────────────────────────
\section{Data}

All data are sourced from RealGM via a custom scraper into a local SQLite
database.  The 2025--26 NCAA Division~I season contains 6{,}305 unique games
involving approximately 350 teams.  Each game yields two observations (one per
team side), for a total of 12{,}610 rows.

The primary outcome variable is offensive efficiency:
\begin{equation}
  e_{\text{off},ig} = \frac{\text{pts}_{ig}}{\text{poss}_{ig}} \times 100
  \label{eq:eoff}
\end{equation}
where $\text{pts}_{ig}$ is points scored and $\text{poss}_{ig}$ is possessions
used by team~$i$ in game~$g$.  Possessions are estimated from boxscore data
using the standard formula.  Observations with fewer than 25 possessions
(garbage time, technical fouls) are excluded.

For the 2025--26 season, the distribution of $e_{\text{off}}$ has mean
$\approx 107$~pts/100 and standard deviation $\approx 16$~pts/100, with a
near-Gaussian shape and heavy right tail from blowout victories.

% ─────────────────────────────────────────────────────────────────────────────
\section{Models}

\subsection{Model 1: Fixed-Point Iteration}

Model~1 follows the KenPom methodology \citep{pomeroy2024}.  Each team $i$ is
assigned latent scalars $(O_i, D_i, R_i) \in \R^3$ representing adjusted
offensive efficiency, adjusted defensive efficiency, and adjusted pace.

\paragraph{Update equations.}  Let $w_g = p_g / \bar{p}$ be possession weights
and $L(h) \in \{1.014, 1.000, 0.986\}$ a home/neutral/away location factor.
The self-consistency conditions are:
\begin{align}
  O_i^{(t+1)} &= \frac{\sum_g w_g \cdot e_{\text{off},ig} \cdot (\bar\mu / D_j^{(t)}) \cdot L(h)^{-1}}
                     {\sum_g w_g}
  \label{eq:O_update} \\
  D_i^{(t+1)} &= \frac{\sum_g w_g \cdot e_{\text{off},jg} \cdot (\bar\mu / O_j^{(t)}) \cdot L(h)}
                     {\sum_g w_g}
  \label{eq:D_update}
\end{align}
where the sums are over games involving team $i$, and $j$ denotes the opposing
team.  Each iteration renormalizes so that $\E_i[O_i] = \E_i[D_i] = \bar\mu$.
Convergence is declared when $\max_i |O_i^{(t+1)} - O_i^{(t)}| < 10^{-6}$.

\paragraph{Prediction.}  Given converged estimates, the predicted offensive
efficiency for team $i$ against opponent $j$ at location $h$ is:
\begin{equation}
  \hat{e}_{\text{off}} = O_i \cdot \frac{D_j}{\bar\mu} \cdot L(h)
  \label{eq:m1_predict}
\end{equation}

\paragraph{Uncertainty.}  Model~1 has no analytic posterior.  We use a
game-level parametric bootstrap: resample games (with both rows of each game
moving together to preserve within-game correlation), refit the fixed-point
iteration, and collect the resulting $(O_i, D_i, R_i)$ as posterior draws.

\paragraph{Limitation.}  Model~1 has no regularization.  With $T \approx 350$
teams and a sparse early-season schedule graph, the system of equations in
\eqref{eq:O_update}--\eqref{eq:D_update} is nearly underdetermined.  Teams
with few games receive extreme estimates that are subsequently corrected only as
more games are played.  This instability is the core motivation for Model~2.

% ─────────────────────────────────────────────────────────────────────────────
\subsection{Model 2: Ridge Regression with Exact Gaussian Posterior}
\label{sec:model2}

\paragraph{Observation model.}  The additive generative model for offensive
efficiency is:
\begin{equation}
  e_{\text{off},ig} = \mu + o_i - d_j + \eta \cdot h_g + \varepsilon_g,
  \qquad \varepsilon_g \sim \N\!\left(0,\, \frac{\sigma^2}{p_g}\right)
  \label{eq:m2_obs}
\end{equation}
where $\mu$ is the league mean, $o_i \in \R$ is team $i$'s offensive
coefficient, $d_j \in \R$ is team $j$'s defensive coefficient (higher = better
defense), $h_g \in \{-1, 0, +1\}$ encodes home/neutral/away for team $i$, and
$\eta$ is the home-court advantage.  Observations are weighted by possessions
$p_g$ to downweight short games.

\paragraph{Prior.}  We place independent Gaussian priors on team effects:
\begin{equation}
  o_i \overset{\mathrm{iid}}{\sim} \N(0, \lambda^{-1}), \qquad
  d_i \overset{\mathrm{iid}}{\sim} \N(0, \lambda^{-1})
  \label{eq:prior}
\end{equation}
with $\lambda_{\text{team}} = 100.0$ (prior variance $\sigma^2 = 0.01$, equivalent to a prior belief that team effects lie within $\pm 0.3$ pts/100 of zero with 95\% probability before seeing data).  No prior is placed on $\mu$ or $\eta$ (i.e.\ $\lambda_\mu = \lambda_\eta = 10^{-4}$, near-flat).

\paragraph{Posterior.}  Let $\theta = (\mu, o_1,\ldots,o_T, d_1,\ldots,d_T, \eta)^\top \in \R^{2T+2}$ and assemble the design matrix $X \in \R^{N \times (2T+2)}$ with rows encoding \eqref{eq:m2_obs}.  Let $W = \mathrm{diag}(p_1/\bar{p}, \ldots, p_N/\bar{p})$ and $\Lambda = \mathrm{diag}(\lambda_\mu, \lambda, \ldots, \lambda, \lambda_\eta)$.  The MAP estimate is:
\begin{equation}
  \hat\theta = \left(X^\top W X + \Lambda\right)^{-1} X^\top W y
  \label{eq:ridge}
\end{equation}
Since the prior is Gaussian and the likelihood is Gaussian, the exact posterior is:
\begin{equation}
  p(\theta \mid \mathcal{D}) = \N\!\left(\hat\theta,\;
    \hat\sigma^2 \left(X^\top W X + \Lambda\right)^{-1}\right)
  \label{eq:posterior}
\end{equation}
where $\hat\sigma^2 = (y - X\hat\theta)^\top W (y - X\hat\theta) / (N - 2T - 2)$.
Posterior samples are drawn via Cholesky decomposition of the covariance.

\paragraph{KenPom summary.}  The three-number summary is derived from the posterior:
\begin{equation}
  \text{AdjO}_i = \mu + o_i, \qquad
  \text{AdjD}_i = \mu - d_i, \qquad
  \text{AdjPace}_i = \exp(\gamma_0 + r_i)
  \label{eq:kenpom}
\end{equation}
where $(\gamma_0, r_1,\ldots,r_T)$ are fit by a separate log-linear pace model.
Crucially, these are \emph{derived} posterior summaries, not primitives: each draw from the posterior yields a full set of $(O_i, D_i, \text{Pace}_i)$ values, so uncertainty propagates through to the KenPom numbers.

\paragraph{Pace model.}  Pace is modeled log-linearly with one row per game:
\begin{equation}
  \log \text{poss}_g = \gamma_0 + r_i + r_j + \gamma_h \cdot |h_g| + \varepsilon_g^{(r)}
\end{equation}
where both teams receive a $+1$ in their respective $r$ columns.

% ─────────────────────────────────────────────────────────────────────────────
\section{Evaluation Protocol}
\label{sec:eval}

\paragraph{One-step-ahead (OSA) forecasting.}
The evaluation protocol is designed to strictly mirror deployment conditions.
Let $\mathcal{G}_t = \{g : \text{date}(g) < t\}$ denote all games played before
date $t$.  For each week $t$ in the season:
\begin{enumerate}
  \item Fit the model on $\mathcal{G}_t$
  \item Predict efficiency for each game $g$ with $t \leq \text{date}(g) < t + 7$
  \item Record residuals $r_{ig} = \hat{e}_{\text{off},ig} - e_{\text{off},ig}$
\end{enumerate}
We require $|\mathcal{G}_t| \geq 300$ before the first prediction window, which begins the second week of November.

This protocol is stricter than the common 80/20 train/test split: the model
never touches any game before predicting it.  The 80/20 split is used only for
calibration analysis in Section~\ref{sec:calibration}.

\paragraph{Metrics.}
For each week $t$ we report:
\begin{align*}
  \text{RMSE}_t &= \sqrt{\frac{1}{N_t}\sum_{g \in W_t}\sum_{i} r_{ig}^2} \\
  \text{MAE}_t  &= \frac{1}{N_t}\sum_{g \in W_t}\sum_i |r_{ig}| \\
  \text{bias}_t &= \frac{1}{N_t}\sum_{g \in W_t}\sum_i r_{ig}
\end{align*}
where $W_t$ is the set of observations in week $t$.  Cumulative versions
aggregate all residuals through week $t$.

% ─────────────────────────────────────────────────────────────────────────────
\section{Results}
\label{sec:results}

\subsection{Early-Season Instability of the Fixed-Point System}

Table~\ref{tab:weekly_rmse} summarizes RMSE for selected weeks.  In week~1
(mid-November, $|\mathcal{G}_t| \approx 700$ games across 350 teams), Model~1
achieves RMSE~$= 22.9$ while Model~2 opens at~$14.6$.

\begin{table}[h]
\centering
\caption{One-step-ahead RMSE (pts/100 poss) for selected weeks, 2025--26 season.}
\label{tab:weekly_rmse}
\begin{tabular}{lrrr}
\toprule
Week & Training games & Model 1 & Model 2 \\
\midrule
Nov 14 &   676 & 22.90 & 14.63 \\
Nov 21 & 1{,}032 & 17.83 & 15.80 \\
Nov 28 & 1{,}425 & 15.14 & 14.29 \\
Dec 05 & 1{,}748 & 14.31 & 13.94 \\
Jan 02 & 2{,}797 & 14.75 & 14.28 \\
Feb 06 & 4{,}503 & 13.79 & 13.68 \\
Mar 13 & 6{,}147 & 12.35 & 12.27 \\
\midrule
\textbf{Season} & 6{,}305 & \textbf{15.10} & \textbf{14.14} \\
\bottomrule
\end{tabular}
\end{table}

Model~1 does not stabilize (RMSE~$<15.5$) until approximately December 5th ---
roughly three weeks into the season.  During those three weeks, Model~2 is
already operating near its season-average accuracy.  We interpret this as the
empirical value of the ridge prior: it encodes the same information that three
weeks of games provide to Model~1.

By mid-January, the gap narrows to $\approx 0.1$--0.2~pts/100.  The cumulative
season RMSE gap is 0.96~pts/100 (14.14 vs.\ 15.10), dominated by the early-season
disparity.

\begin{proposition}[Informal]
  The ridge prior in Model~2 is equivalent in forecast value to approximately
  three additional weeks of game data for Model~1 in early-season prediction.
\end{proposition}

This is not a theoretical claim but an empirical observation: Figure~\ref{fig:rmse_curves}
(see notebook) shows cumulative RMSE curves converging, with Model~2's initial
level matching Model~1's level at the point of stabilization $\approx$~December~5th.

\subsection{Season Summary Statistics}
\label{sec:season_stats}

Table~\ref{tab:season_summary} gives season-wide metrics across 10{,}562
test observations.

\begin{table}[h]
\centering
\caption{Season-wide one-step-ahead forecast metrics, 2025--26.}
\label{tab:season_summary}
\begin{tabular}{lrrr}
\toprule
Model & RMSE & MAE & Bias \\
\midrule
Model 1 (fixed-point) & 15.097 & 11.909 & $-$1.449 \\
Model 2 (ridge)       & 14.138 & 11.244 & $-$1.175 \\
\bottomrule
\end{tabular}
\end{table}

% ─────────────────────────────────────────────────────────────────────────────
\subsection{Calibration of Posterior Predictive Intervals}
\label{sec:calibration}

To assess calibration, we fit Model~2 on the first 80\% of games (by game
identifier, approximately chronological) and construct posterior predictive
intervals for the held-out 20\%.

The predictive distribution for a single observation is:
\begin{equation}
  p(\tilde{e}_{ij} \mid \mathcal{D}) \approx \frac{1}{S}\sum_{s=1}^S
  \N\!\left(\hat{e}_{ij}^{(s)},\; \sigma^2_{\text{eff}}\right),
  \qquad \theta^{(s)} \sim p(\theta \mid \mathcal{D})
  \label{eq:pp_mixture}
\end{equation}
where $\sigma^2_{\text{eff}} = \E[(\hat{e} - e_{\text{off}})^2]$ is estimated
from training residuals.

The calibration curve (Figure~\ref{fig:calibration}) plots empirical coverage
against nominal coverage for 60 levels from 2\% to 99\%.  Reference values:

\begin{center}
\begin{tabular}{cc}
\toprule
Nominal & Empirical \\
\midrule
50\% & 49\% \\
80\% & 80\% \\
90\% & 88\% \\
95\% & 93\% \\
\bottomrule
\end{tabular}
\end{center}

The model is \emph{very slightly underconfident} at high coverage levels (95\% nominal $\to$ 93\% empirical), meaning the intervals are marginally wider than necessary.  This is the safer failure mode: overconfident intervals (too narrow) miss actual outcomes and mislead downstream decisions.

\subsection{The Bias Signal}
\label{sec:bias}

Both models exhibit a consistent negative bias of $\approx -1.2$ to $-1.5$~pts/100
(predictions are below actuals).  This pattern is not attributable to model
misspecification in the standard sense.

The model assumes $o_i$ is constant within the season.  In reality, team
offenses improve as players gain experience, coaching adjustments take hold, and
injured players return.  A model trained on November and December games and
asked to predict March games will therefore systematically under-predict,
because it has learned an earlier, weaker version of each team's offense.

This interpretation generates a testable prediction: the bias should be larger
for games later in the season.  We observe this in Figure~\ref{fig:bias_curves}
(notebook): the weekly bias is near zero in early December and becomes
consistently negative by late January, growing in magnitude through the
tournament.

The natural fix is to replace the static $o_i$ with a dynamic state:
\begin{equation}
  o_i(t) = o_i(t-1) + \xi_t, \qquad \xi_t \sim \N(0, \tau^2)
  \label{eq:dynamic}
\end{equation}
yielding a Kalman-filter-style inference problem.  We leave this extension for
future work.

% ─────────────────────────────────────────────────────────────────────────────
\section{Team Rating Trajectories and Uncertainty}

Re-fitting Model~2 weekly on cumulative data produces team net-rating
trajectories $\hat{o}_i(t) + \hat{d}_i(t)$ with posterior $\pm 1\sigma$ and
$\pm 2\sigma$ bands at each time point.

Several observations follow from the 2025--26 season:

\begin{itemize}
  \item Early-season ($|\mathcal{G}| < 1500$): posterior $\sigma$ for a typical
    top-25 team is $1.5$--$3.0$~pts/100.  At this scale, most teams' credible
    intervals overlap substantially.
  \item By mid-January ($|\mathcal{G}| \approx 3800$): $\sigma$ has collapsed to
    $0.3$--$0.4$~pts/100 for teams with 15+ games.  The hierarchy is clear.
  \item The rate of uncertainty reduction is approximately uniform across teams
    at similar stages of the schedule, confirming that variance is driven by
    sample size rather than structural instability.
  \item Elite teams (Michigan, Duke, Arizona in 2025--26) achieve smaller
    late-season $\sigma$ than mid-major teams, consistent with playing more
    high-information (competitive) games.
\end{itemize}

% ─────────────────────────────────────────────────────────────────────────────
\section{Pre-Game Predictive Distributions}

For any forthcoming game between teams $i$ and $j$, the full predictive
distribution \eqref{eq:pp_mixture} is available for both teams.  The overlap
between $p(\tilde{e}_i \mid \mathcal{D})$ and $p(\tilde{e}_j \mid \mathcal{D})$ provides a
visual and quantitative measure of competitiveness.

For six top-25 matchups sampled across the 2025--26 season, we observe:
\begin{itemize}
  \item Games between teams with similar net ratings produce nearly identical
    predictive densities, reflecting genuine uncertainty about outcomes.
  \item Games involving a large net-rating gap show clearly separated
    distributions, with actual outcomes (dotted lines) falling within the bulk
    of the predicted density in all six cases --- a direct visual calibration check.
  \item The width of individual densities ($\approx 25$--30~pts/100 at the 95\% PI)
    is driven primarily by $\sigma_{\text{eff}} \approx 14$~pts/100 (irreducible
    game randomness), with parameter uncertainty contributing $\approx 1$--2~pts/100.
\end{itemize}

% ─────────────────────────────────────────────────────────────────────────────
\section{Discussion}

\subsection{Ratings as Posterior Summaries}

The central conceptual contribution of this work is framing team ratings as
\emph{derived summaries of a posterior distribution}, rather than primitives.
Formally:
\begin{equation}
  \theta \sim p(\theta \mid \mathcal{D})
  \quad\Longrightarrow\quad
  S(\theta) = (\text{AdjO}_i, \text{AdjD}_i, \text{AdjPace}_i)_{i=1}^T
\end{equation}

This means that every downstream quantity --- KenPom-style ratings, matchup
predictions, win probabilities --- inherits calibrated uncertainty from the same
generative object.  The framework subsumes the point-estimate case as the
posterior mean, while providing uncertainty estimates that are provably
calibrated.

\subsection{Comparison to Existing Work}

Bayesian approaches to sports ratings have a long history
\citep{glickman1999,kovalchik2016}.  The closest methodological relative is the
RAPM (Regularized Adjusted Plus-Minus) framework from basketball analytics
\citep{sill2010}, which uses ridge regression on player-level data.  Our
contribution extends this to team-level ratings with an explicit posterior and
a rigorous evaluation protocol.

The key novelty is the one-step-ahead validation, which is largely absent in
the sports analytics literature.  Most published evaluations use either a
single train/test split or report in-sample fit, neither of which reflects
deployment conditions.

\subsection{Limitations and Extensions}

\paragraph{Static coefficients.}  The primary limitation is the assumption of
time-invariant $o_i$ and $d_i$.  The consistent negative bias (Section~\ref{sec:bias}) is a direct consequence.  The Kalman filter extension in \eqref{eq:dynamic} is the natural remedy.

\paragraph{Schedule strength circularity.}  Adjusted efficiency is computed
relative to opponents' adjusted efficiency, creating a circular dependency that
the fixed-point iteration resolves for Model~1 and the ridge regression
approximates for Model~2 (via the $o_i - d_j$ encoding).  True opponent-adjusted
inference would require a hierarchical model.

\paragraph{Player-level structure.}  Team-level ratings aggregate over lineup
combinations.  The next modeling step is a player-level decomposition where
team efficiency is a function of the active lineup, enabling transfer of
information across teams sharing personnel.  This connects naturally to the
graph neural network architecture of \citet{danielson2024deeprapm}.

% ─────────────────────────────────────────────────────────────────────────────
\section{Conclusion}

We have demonstrated that a ridge prior over team effects in a Gaussian
observation model produces:
\begin{enumerate}
  \item Early-season forecast accuracy equivalent to three additional weeks of
    game data compared to the standard fixed-point approach.
  \item Season-wide RMSE of 14.14~pts/100 versus 15.10 for the fixed-point
    system, under strict one-step-ahead evaluation.
  \item Empirically near-perfect calibration across all nominal coverage levels,
    with slight and benign underconfidence at high levels.
  \item Interpretable systematic bias as a within-season team improvement
    signal, pointing to a dynamic-coefficient extension.
\end{enumerate}

The system produces a full posterior distribution over team quality, enabling
calibrated probabilistic statements about matchup outcomes.  Ratings, predictions,
and win probabilities all flow from the same generative object with
quantifiable uncertainty at every level.

% ─────────────────────────────────────────────────────────────────────────────
\section*{Reproducibility}

All code is available in the accompanying repository.  The dataset is derived
from RealGM via an open scraper.  The notebook
\texttt{notebooks/probabilistic\_ratings.ipynb} reproduces all figures and
tables in this paper.

% ─────────────────────────────────────────────────────────────────────────────
\bibliographystyle{plainnat}
\begin{thebibliography}{9}

\bibitem[Danielson(2024)]{danielson2024deeprapm}
Danielson, A.\ (2024).
\newblock Deep RAPM: Player embeddings via set transformers.
\newblock \emph{Unpublished manuscript}.

\bibitem[Glickman and Stern(1998)]{glickman1999}
Glickman, M.~E.\ and Stern, H.~S.\ (1998).
\newblock A state-space model for national football league scores.
\newblock \emph{Journal of the American Statistical Association},
  93(441):25--35.

\bibitem[Kovalchik(2016)]{kovalchik2016}
Kovalchik, S.~A.\ (2016).
\newblock Searching for the GOAT of tennis win prediction.
\newblock \emph{Journal of Quantitative Analysis in Sports}, 12(3):127--138.

\bibitem[Pomeroy(2024)]{pomeroy2024}
Pomeroy, K.\ (2024).
\newblock Ratings explanation.
\newblock \url{https://kenpom.com/blog/ratings-explanation/}.

\bibitem[Sill(2010)]{sill2010}
Sill, J.\ (2010).
\newblock Improved NBA adjusted plus-minus using regularization and
  out-of-sample testing.
\newblock In \emph{Proceedings of the 2010 MIT Sloan Sports Analytics
  Conference}.

\end{thebibliography}

\end{document}
